\clearpage
\fancyhead[L]{\small CAGNOLI. STERNVERZEICHNISS.}
\hypertarget{native-501-lower}{}\label{native-501-lower}
\begin{center}
\small Göttingische gelehrte Anzeigen. No. 77. pp. 761..763. 14 May 1808.\\[3pt]
\rule{.72\linewidth}{.3pt}
\end{center}

\emph{Catalogue de \textup{501} étoiles suivi des tables relatives d'Aberration et de Nutation, par Antoine Cagnoli.} Modena 1807. 280 pages in quarto.

This star catalogue had already appeared in volume X of the \emph{Memorie di Matematica e di Fisica della Societa Italiana} and is published here in an improved and enlarged form. That first printing was also accompanied by historical information about how the catalogue came to be, from which we select a few details here. \textsc{Cagnoli} began compiling a catalogue of stars from his own observations as early as 1783 in Paris, where he was then living, and subsequently completed the work in Verona in the years 1788\mbox{. . .}1792 with the same instruments. His plan was to give at least four reliable determinations of stars in each zone north of the equator spanning half a degree in declination, with the stars about 6 hours apart in right ascension. This explains the selection of stars, which is striking at first sight. The stars chosen are by no means the brightest in every constellation. In their place the catalogue often includes stars down to the sixth magnitude. \textsc{Cagnoli}, however, could not adhere to his plan entirely, partly because suitable stars could not be found everywhere, and partly because his health and other interruptions prevented him. \textsc{Cagnoli} wanted to borrow nothing from others for his catalogue. To obtain absolute right ascensions, he therefore determined that of $\alpha$ in Auriga by comparing it directly with the Sun 24 times, and found a result that agreed with \textsc{Maskelyne}’s latest value to within $1''$. For the entire work he used a movable quadrant of three feet in radius, with an achromatic telescope with a two-inch aperture, and a 3$\frac{1}{2}$-foot achromatic meridian telescope with an aperture of 28 lines, both by \textsc{Megnié}, and an excellent pendulum clock made by \textsc{Robins} in Paris.
\SourceLeaf{502}
All these instruments are now at the observatory in Milan. Some stars, incidentally, were determined from observations by \textsc{Cesaris} in Milan and are distinguished in the catalogue by a special mark.

Most of the 501 stars in the catalogue are northern, and only 28 are southern. Their positions are given for 1800. Each right ascension and declination is accompanied by the number of observations and the difference between those that deviate most from one another. For most stars the differences from \textsc{Piazzi}’s determinations amount to only a few seconds. It really is evidence of exceptional skill and care that \textsc{Cagnoli} was able to achieve this degree of precision with equipment that was only moderately good for such a purpose. At the end of the catalogue all the stars are arranged again by their distances from the north pole, which is very convenient for anyone who wants to use them for comparison.

Of particular value are the special aberration and nutation tables for most of these stars, included with the new printing of the catalogue. The aberration tables are already calculated using \textsc{Delambre}’s latest determination of the constant of aberration ($20''25$). The nutation tables, on the other hand, still assume $9''0$ and $6''7$ for the major and minor semiaxes of the nutation ellipse. All the numbers in these tables must therefore be increased by one ninth or one fourteenth if they are to be brought into agreement with the latest values of \textsc{Laplace} or \textsc{von Zach}.
